% demo-myinclude.tex -- \myinclude
\documentclass[sectioning=nobreak]{myreport}
\usepackage{mydemo}
\begin{document}
\demoheader{\string\myinclude}{Includes an image (not a float, no caption),
sized relative to the text block.  Used by \string\myfigure,
\string\mysubfigure, \string\mylogo\ and \string\myimage.}

\mysection*{Syntax}
\begin{demosyntax}
\myinclude[<options>]{<file>}
\end{demosyntax}
\begin{itemize}
  \item \texttt{height=<f>}: height = $f\times$\verb|\textheight|;
        \texttt{width=<f>}: width = $f\times$\verb|\textwidth| (with
        the default base \texttt{text}). A length (\texttt{3cm}) is used
        as it is. \texttt{base=}: \texttt{text}, \texttt{none} (a unit
        is required) or a unit (\texttt{base=cm, width=3} is 3\,cm).
        \texttt{natural}: natural size. The default,
        \texttt{height=0.2}, is set by the package option \texttt{size}.
        With both keys the image keeps its proportions and fits in the
        box.
  \item \texttt{<file>}: name with extension, in the \texttt{images}
        folder. \texttt{.svg} files are converted with Inkscape
        (\texttt{-shell-escape} needed).
\end{itemize}

\mysection*{Examples}
\begin{demo}
\myinclude{wide.png}
\end{demo}
\begin{demo}
\myinclude[width=0.3]{wide.png}\hfill\myinclude[height=0.1]{tall.png}
\end{demo}
\begin{demo}
\myinclude[height=0.05, width=0.5]{square.jpg} % fits: 0.05 high
\hfill \myinclude[base=cm, width=3]{wide.png}   % 3 cm
\hfill \myinclude[width=2cm]{wide.png}          % 2 cm
\end{demo}
\begin{demo}
In a sentence: \myinclude[height=0.02]{square.jpg} inline.
\end{demo}
\begin{demo}
\myinclude[width=0.4]{plot.svg}
\end{demo}
\end{document}
